\documentclass[11pt,letterpaper]{article}

\usepackage[letterpaper,top=0.45in,bottom=0.5in,left=0.5in,right=0.5in]{geometry}
\usepackage{graphicx}
\usepackage{xcolor}
\usepackage{listings}
\usepackage{booktabs}
\usepackage{hyperref}
\usepackage{caption}
\usepackage{float}
\usepackage{titlesec}
\usepackage{amsmath}

\hypersetup{
    colorlinks=true,
    linkcolor=blue!60!black,
    urlcolor=blue!60!black,
}

\setlength{\parskip}{2pt}
\setlength{\parindent}{0pt}
\titlespacing*{\section}{0pt}{6pt}{2pt}
\titlespacing*{\subsection}{0pt}{4pt}{2pt}
\captionsetup{skip=2pt,font=small}

\graphicspath{{figures/}}

% -----------------------------------------------------------------------------
% Optional figures: if the screenshot has not yet been captured from
% Vivado/Vitis, fall back to a labelled framed box so the rest of the
% report still compiles cleanly.
% -----------------------------------------------------------------------------
\usepackage{tikz}
\newcommand{\optionalfig}[3][\textwidth]{%
    \IfFileExists{figures/#2}{%
        \includegraphics[width=#1]{#2}%
    }{%
        \begin{tikzpicture}
            \draw[dashed,thick] (0,0) rectangle (#1,4cm);
            \node[align=center] at (0.5*#1,2cm) {%
                \textit{placeholder:} \texttt{#2}\\[2pt]
                {\footnotesize #3}%
            };
        \end{tikzpicture}%
    }%
}

% -----------------------------------------------------------------------------
% Code listing styles
% -----------------------------------------------------------------------------
\definecolor{codebg}{RGB}{248,248,248}
\definecolor{codekw}{RGB}{0,0,180}
\definecolor{codecmt}{RGB}{0,120,0}
\definecolor{codestr}{RGB}{170,55,20}
\definecolor{codenum}{RGB}{120,120,120}

\lstset{
    language=C,
    basicstyle=\ttfamily\footnotesize,
    keywordstyle=\color{codekw}\bfseries,
    commentstyle=\color{codecmt}\itshape,
    stringstyle=\color{codestr},
    numbers=left,
    numberstyle=\tiny\color{codenum},
    stepnumber=1,
    numbersep=5pt,
    backgroundcolor=\color{codebg},
    frame=single,
    rulecolor=\color{black!30},
    tabsize=4,
    showstringspaces=false,
    breaklines=true,
    breakatwhitespace=false,
    breakindent=0pt,
    postbreak=\mbox{\textcolor{red}{$\hookrightarrow$}\space},
    captionpos=b,
    aboveskip=4pt,
    belowskip=4pt,
    xleftmargin=6pt,
    xrightmargin=2pt,
    columns=fullflexible,
    keepspaces=true,
    morekeywords={u32,s32,Xil_Out32,Xil_In32}
}

\lstdefinelanguage{VHDLsmall}{
    sensitive=false,
    morekeywords={architecture,begin,case,component,downto,else,elsif,end,
        entity,if,is,library,others,port,map,signal,std_logic,
        std_logic_vector,then,to,use,when,work,all,generic,return,process,
        rising_edge,unsigned,signed,std_logic_unsigned,numeric_std,
        constant,variable,type,array,natural,integer,for,loop,while},
    morecomment=[l]--
}

\lstdefinestyle{vhdlstyle}{
    language=VHDLsmall,
    basicstyle=\ttfamily\footnotesize,
    keywordstyle=\color{codekw}\bfseries,
    commentstyle=\color{codecmt}\itshape,
    numbers=left,
    numberstyle=\tiny\color{codenum},
    numbersep=5pt,
    backgroundcolor=\color{codebg},
    frame=single,
    rulecolor=\color{black!30},
    tabsize=4,
    breaklines=true,
    captionpos=b,
    aboveskip=4pt,
    belowskip=4pt,
    xleftmargin=6pt,
    xrightmargin=2pt,
    columns=fullflexible,
    keepspaces=true
}

\title{\vspace{-0.6in}\textbf{ECEC 661 -- Homework 5}\\
       \large Stack Processor 101 with new \texttt{scpb} block-copy instruction}
\author{Gary Pham \texttt{(gp492)}}
\date{\today}

\begin{document}
\maketitle
\vspace{-0.25in}

% =============================================================================
\section{Problem Statement}
% =============================================================================

Starting from the assignment PDF's Stack Processor~101 (\texttt{sp101}),
implement a new machine instruction \texttt{scpb} -- \emph{stack copy
block} -- which pops three arguments from the data stack and copies a
contiguous block of memory. The user-visible calling convention is:
\[
\texttt{sc}\ \texttt{dest\_addr}\quad
\texttt{sc}\ \texttt{source\_addr}\quad
\texttt{sc}\ \texttt{number\_of\_data}\quad
\texttt{scpb}
\]
At entry the stack top holds \texttt{number\_of\_data}, then
\texttt{source\_addr}, then \texttt{dest\_addr}. \texttt{scpb} pops in
that order and copies \texttt{number\_of\_data} consecutive 32-bit words
from \texttt{source\_addr} into the contiguous block starting at
\texttt{dest\_addr}. Deliverables: the updated user logic, behavioural
simulation evidence of correctness, IP packaging of \texttt{sp101}, and
a Vitis test application with proof on the Cora~Z7-07S.

% =============================================================================
\section{Architecture}
% =============================================================================

The IP keeps the assignment PDF's three-block layout
(Figure~\ref{fig:bd}): an AXI4-Lite slave exposes the five
software-visible registers, a \texttt{bus\_ip\_mem\_bridge} multiplexes
the bus or the processor onto the single port of
\texttt{blk\_mem\_gen\_0}, and the stack processor (\texttt{user\_logic})
fetches/executes machine instructions out of that BRAM.

\begin{figure}[H]
    \centering
    \optionalfig[0.92\textwidth]{diagram.png}{Vivado block design screenshot
        (Zynq PS, AXI interconnect, \texttt{sp101\_axi\_0}).}
    \caption{Vivado block design: Zynq PS, AXI interconnect, and the custom
             \texttt{sp101\_axi} IP (\texttt{user\_logic} + bridge + BRAM
             inside).}
    \label{fig:bd}
\end{figure}

% =============================================================================
\section{Software Register Map}
% =============================================================================

The slave maps five 32-bit registers into a 32-byte aperture, identical
to the PDF's reference C snippet:

\begin{table}[H]
    \centering
    \caption{Software register map (byte offsets from
             \texttt{XPAR\_SP101\_AXI\_0\_S00\_AXI\_BASEADDR}).}
    \label{tab:regs}
    \small
    \begin{tabular}{clcp{8.0cm}}
        \toprule
        \textbf{Offset} & \textbf{Name} & \textbf{R/W} & \textbf{Purpose} \\
        \midrule
        \texttt{0x00} & \texttt{slv\_reg0} & R/W &
            control: bit~0~=\texttt{run}, bit~1~=\texttt{reset},
            bit~2~=\texttt{bus2mem\_en}, bit~3~=\texttt{bus2mem\_we} \\
        \texttt{0x04} & \texttt{slv\_reg1} & R/W &
            \texttt{bus2mem\_addr} (lower 10 bits) \\
        \texttt{0x08} & \texttt{slv\_reg2} & R/W &
            \texttt{bus2mem\_data\_in} (32 bit) \\
        \texttt{0x0C} & \texttt{slv\_reg3} & R   &
            \texttt{sp2bus\_data\_out} (32 bit) \\
        \texttt{0x10} & \texttt{slv\_reg4} & R   &
            bit~0 = \texttt{done} flag \\
        \bottomrule
    \end{tabular}
\end{table}

% =============================================================================
\section{Instruction Set and \texttt{scpb} Opcode}
% =============================================================================

The processor honours every opcode from the assignment PDF and adds a
new family for \texttt{scpb}. The lowest hex digit of each constant is
the micro-step index inside the instruction; the higher hex digits are
the opcode group.

\begin{table}[H]
    \centering
    \caption{Instruction set (opcode entry constants in
             \href{run:../stackprocessor_101/rtl/user_logic.vhd}{\texttt{user\_logic.vhd}}).}
    \label{tab:isa}
    \small
    \begin{tabular}{llll}
        \toprule
        \textbf{Mnemonic} & \textbf{Code} & \textbf{Stack effect} & \textbf{Notes} \\
        \midrule
        \texttt{halt} & \texttt{0x000000FF} & --                                            & raise \texttt{done}, halt \\
        \texttt{sc}   & \texttt{0x00000001} & push \(c\)                                    & push constant pointed by pc \\
        \texttt{sl}   & \texttt{0x00000011} & pop \(a\), push \(\mathrm{mem}[a]\)           & indirect load \\
        \texttt{ss}   & \texttt{0x00000021} & pop \(d\), pop \(a\), \(\mathrm{mem}[a]\)~\(\leftarrow d\) & indirect store \\
        \texttt{sadd} & \texttt{0x00000031} & pop, pop, push sum                            & stack add \\
        \texttt{scp}  & \texttt{0x00000101} & pop \(s\), pop \(d\), \(\mathrm{mem}[d]\)~\(\leftarrow \mathrm{mem}[s]\) & single-word copy \\
        \texttt{scpb} & \texttt{0x00000201} & pop \(N\), pop \(s\), pop \(d\), block copy   & \textbf{NEW} \\
        \bottomrule
    \end{tabular}
\end{table}

\subsection{\texttt{scpb} micro-program}

The FSM (\texttt{Homework\_5/stackprocessor\_101/rtl/user\_logic.vhd},
state \texttt{exe} branch) walks \texttt{scpb} through three phases:
pop the three arguments, loop over the copy, and per-iteration
bookkeeping. The PDF's simulation latency cushion is preserved: each
BRAM read step has the same number of extra ``mem\_addr settled /
douta valid / mem\_data\_out latched'' micro-states the legacy opcodes
use.

\begin{lstlisting}[style=vhdlstyle,caption={Core micro-states (excerpt of
\texttt{when SCPB ...} body in \texttt{user\_logic.vhd}).},label={lst:scpb}]
when SCPB =>                            -- init pop count
    mem_addr <= std_logic_vector(unsigned(sp) - 1);
    sp       <= std_logic_vector(unsigned(sp) - 1);
    we       <= "0"; ir <= SCPB2;
...
when SCPB5 =>                           -- count valid -> copy_cnt
    copy_cnt <= mem_data_out;
    ir       <= SCPB6;
when SCPB6 =>                           -- source valid -> copy_src
    copy_src <= mem_data_out(A-1 downto 0);
    ir       <= SCPB7;
when SCPB7 =>                           -- dest valid -> copy_dst, enter loop
    copy_dst <= mem_data_out(A-1 downto 0);
    ir       <= SCPB_L;
when SCPB_L =>
    if unsigned(copy_cnt) = 0 then
        n_s <= fetch;                   -- block copy finished
    else
        mem_addr <= copy_src; ir <= SCPB_R1;
    end if;
when SCPB_R4 =>                         -- source word -> schedule write
    mem_addr    <= copy_dst;
    mem_data_in <= mem_data_out;
    we          <= "1"; ir <= SCPB_W;
when SCPB_W =>                          -- write fires, bookkeeping
    copy_src <= std_logic_vector(unsigned(copy_src) + 1);
    copy_dst <= std_logic_vector(unsigned(copy_dst) + 1);
    copy_cnt <= std_logic_vector(unsigned(copy_cnt) - 1);
    we       <= "0"; ir <= SCPB_L;
\end{lstlisting}

Pre-loop overhead is 7 cycles, each iteration costs 6 cycles, and the
terminal check adds one more, so a block of $N$ words takes roughly
$6N+8$ clocks plus the usual \texttt{fetch} overhead.

% =============================================================================
\section{Behavioral Simulation}
% =============================================================================

\texttt{tb/tb\_user\_logic.vhd} drives the same \texttt{run},
\texttt{reset}, \texttt{bus2mem\_*} signals the assignment's xsim script
forces, walking three self-checking scenarios:

\begin{enumerate}
    \item \emph{multi-word}: copy three words
          (\texttt{0x0B,0x16,0x21}) from addresses \(100\ldots102\) to
          \(200\ldots202\),
    \item \emph{single word}: copy one word (\texttt{0x63}) from
          address \(110\) to \(210\) (off-by-one guard),
    \item \emph{zero count}: pre-fill \(220 = \mathtt{0xDEADBEEF}\), run
          \texttt{scpb} with \texttt{count}=0, then verify the
          destination is unchanged (no-op guard).
\end{enumerate}

Each scenario emits one \texttt{[PASS]} per destination word followed by
the testbench banner. Figure~\ref{fig:sim} shows the full xsim run
spanning all three scenarios; the integer
\texttt{checks} counter climbs from~0 to~5 while \texttt{errors} stays
at~0, confirming every destination word matches its golden value.

\begin{figure}[H]
    \centering
    \optionalfig[\textwidth]{sim.png}{xsim waveform spanning the entire
        \texttt{tb\_user\_logic} run, with \texttt{checks=5} and
        \texttt{errors=0}.}
    \caption{xsim waveform across the full testbench run.  \texttt{n\_s}
             cycles \texttt{idle\,$\to$\,exe\,$\to$\,chill} three times,
             once per scenario.  The block-copy registers walk through
             scenario~1 (\texttt{copy\_src}: 100,101,102~--
             \texttt{copy\_dst}: 200,201,202 -- \texttt{copy\_cnt}:
             3,2,1,0), scenario~2 (110~$\to$~210, count=1), and
             scenario~3 (120/220, count=0 no-op).  The
             \texttt{checks}/\texttt{errors} integers in the wave window
             prove the self-check passes (5/0).}
    \label{fig:sim}
\end{figure}

% =============================================================================
\section{Software (\texttt{main.c})}
% =============================================================================

The bare-metal application
(\texttt{Homework\_5/stackprocessor\_101/sw/main.c}) drives the AXI4-Lite
peripheral through the exact handshake described in the PDF. It also
follows the same scenario list as the testbench so the on-board UART log
matches the simulator output line-for-line.

\begin{lstlisting}[caption={Core of \texttt{main.c}: write a word to BRAM,
                           run, read it back.},label={lst:handshake}]
/* IMPORTANT: never assert CTRL_RESET while reading.  The mem_data_out
 * register chain inside user_logic is forced to 0 while reset='1', so
 * a read with reset held returns 0x00000000 and silently masks both the
 * pre-written sentinel and the scpb result. */

static void mem_write(u32 addr, u32 data)
{
    SP101_mWriteReg(BASEADDR, REG_CTRL,     CTRL_BUS_EN | CTRL_BUS_WE);
    SP101_mWriteReg(BASEADDR, REG_BUS_ADDR, addr & 0x3FFU);
    SP101_mWriteReg(BASEADDR, REG_BUS_DIN,  data);
}

static u32 mem_read(u32 addr)
{
    SP101_mWriteReg(BASEADDR, REG_CTRL,     CTRL_BUS_EN);
    SP101_mWriteReg(BASEADDR, REG_BUS_ADDR, addr & 0x3FFU);
    for (volatile int s = 0; s < 64; s++) { }
    return SP101_mReadReg(BASEADDR, REG_SP_DOUT);
}

static void processor_reset(void)
{
    SP101_mWriteReg(BASEADDR, REG_CTRL, CTRL_RESET | CTRL_BUS_EN);
    for (volatile int s = 0; s < 64; s++) { }
    SP101_mWriteReg(BASEADDR, REG_CTRL, CTRL_BUS_EN);
}

static void processor_run_and_wait(void)
{
    SP101_mWriteReg(BASEADDR, REG_CTRL, CTRL_RUN);
    while ((SP101_mReadReg(BASEADDR, REG_DONE) & 0x1U) == 0U) { }
    SP101_mWriteReg(BASEADDR, REG_CTRL, 0U);
}
\end{lstlisting}

% =============================================================================
\section{Hardware Results}
% =============================================================================

The same three scenarios were re-executed on the Cora~Z7-07S over the
AXI4-Lite interface.  The UART transcript (Figure~\ref{fig:vitis})
reports every dest word readback as \texttt{OK} and ends with
\texttt{Summary : 3/3 scenarios PASSED}, matching the simulation
result line-for-line.  The exact captured trace is reproduced in
Listing~\ref{lst:uart} for reference.

\begin{figure}[H]
    \centering
    \optionalfig[0.55\textwidth]{vitis.png}{PuTTY / serial-monitor capture
        of the three \texttt{scpb} scenarios on Cora Z7-07S.}
    \caption{UART trace from the bare-metal \texttt{sp101 + scpb} test app
             running on the Cora~Z7-07S: the multi-word
             ($100..102\to200..202$), single-word ($110\to210$), and
             zero-count ($\mathtt{0xDEADBEEF}$ sentinel preserved at
             address~220) scenarios all PASS.}
    \label{fig:vitis}
\end{figure}

\begin{lstlisting}[language={},caption={Captured UART trace
                           (verbatim from the Cora~Z7-07S
                           run).},label={lst:uart}]
==========================================================
   ECEC 661 HW5 - Stack Processor 101 + scpb block copy
   IP base address : 0x43C30000
==========================================================

---- Scenario: scpb count=3 (multi-word) ----
  dest=200  src=100  count=3
  -- destination dump --
    [200] got=0x0000000B  expect=0x0000000B   OK
    [201] got=0x00000016  expect=0x00000016   OK
    [202] got=0x00000021  expect=0x00000021   OK
  RESULT: PASS

---- Scenario: scpb count=1 (single word) ----
  dest=210  src=110  count=1
  -- destination dump --
    [210] got=0x00000063  expect=0x00000063   OK
  RESULT: PASS

---- Scenario: scpb count=0 (no-op) ----
  dest=220  src=120  count=0
  -- destination dump --
    [220] got=0xDEADBEEF  expect=0xDEADBEEF   OK
  RESULT: PASS

----------------------------------------------------------
   Summary : 3/3 scenarios PASSED
----------------------------------------------------------
\end{lstlisting}

\begin{table}[H]
    \centering
    \caption{Reference test vectors for \texttt{scpb} (identical in
             simulation and on hardware).}
    \label{tab:vectors}
    \small
    \begin{tabular}{clllll}
        \toprule
        \textbf{Scenario} & \texttt{count} & \texttt{source} & \texttt{dest} & \textbf{Source words} & \textbf{Expected at dest} \\
        \midrule
        multi-word  & 3 & 100 & 200 & \texttt{0x0B, 0x16, 0x21} & \texttt{0x0B, 0x16, 0x21} \\
        single-word & 1 & 110 & 210 & \texttt{0x63}              & \texttt{0x63}              \\
        zero-count  & 0 & 120 & 220 & \emph{(unused)}            & \texttt{0xDEADBEEF} (unchanged) \\
        \bottomrule
    \end{tabular}
\end{table}

\subsection{Driver gotcha worth flagging}

While bringing the bare-metal app up, all three on-board scenarios
initially failed with \texttt{got=0x00000000} for every dump line
(including the \texttt{0xDEADBEEF} sentinel of the no-op test, which
should never be touched).  The cause was that the original
\texttt{mem\_read()} held \texttt{CTRL\_RESET} while reading
\texttt{slv\_reg3}; the \texttt{mem\_data\_out} register chain in
\texttt{user\_logic.vhd} clears on \texttt{reset='1'}, so every read
returned~0.  The fix shown in Listing~\ref{lst:handshake} only pulses
reset inside \texttt{processor\_reset()} and keeps it low for the rest
of each scenario.  As a defensive belt-and-braces, the
\texttt{p\_mdo} register chain in
\href{run:../stackprocessor_101/rtl/user_logic.vhd}{\texttt{user\_logic.vhd}}
no longer zero-clears on reset, so a future driver misuse cannot
silently mask read data.

% =============================================================================
\section{Conclusion}
% =============================================================================

The Homework~5 design starts from the assignment PDF's Stack
Processor~101, adds a 13-state \texttt{scpb} micro-program for arbitrary
block copies, packages the result as the AXI4-Lite IP \texttt{sp101\_axi},
and exercises every code path (multi-word, single-word, zero-count) in
both behavioural simulation (Figure~\ref{fig:sim}, \texttt{checks=5},
\texttt{errors=0}) and on the Cora~Z7-07S
(Figure~\ref{fig:vitis} / Listing~\ref{lst:uart},
\texttt{3/3 scenarios PASSED}).  Each deliverable
listed in \href{run:../prob.md}{\texttt{Homework\_5/prob.md}} maps to a
concrete artifact under
\href{run:../stackprocessor_101/}{\texttt{Homework\_5/stackprocessor\_101/}}
and is reproduced by following \texttt{README.md}.

\end{document}
